\documentclass[leqno, a4paper, 12pt]{jsarticle}
\usepackage{amssymb}
\usepackage{amsthm}
\usepackage{amsmath}
\usepackage{comment}
\usepackage{url}

\usepackage{enumitem}
%\usepackage{showkeys}


%定理環境 thmはあまり使わずpropをなるべく使う。
%主張は基本的に証明の中で使い、そのため番号を付けない。
%注意環境を付けてみた。
\newtheorem{thm}{定理}
\newtheorem{prop}[thm]{命題}
\newtheorem{cor}[thm]{系}
\newtheorem{lem}[thm]{補題}
\newtheorem{df}[thm]{定義}
\newtheorem{exam}[thm]{例}
\newtheorem{fact}[thm]{事実}
\newtheorem*{claim}{主張}
\newtheorem*{cau}{注意}
\newtheorem*{rmk}{注意}
\renewcommand{\proofname}{証明}
%矢印たち。
%\toは\rightarrowと同じ。\getsは\leftarrowと同じ。同値は\iff
\newcommand{\To}{\Longrightarrow}
\newcommand{\Gets}{\LongLeftarrow}
%黒板太字
\newcommand{\nn}{\mathbb{N}}
\newcommand{\zz}{\mathbb{Z}}
\newcommand{\qq}{\mathbb{Q}}
\newcommand{\rr}{\mathbb{R}}
\newcommand{\cc}{\mathbb{C}}
%無限大
\newcommand{\mgn}{\infty}
%差集合は\setminusで統一する。等号の否定は\neq
%真部分集合は\subsetneqq　偏微分は\partial
%ディスプレイスタイル。\dfracでディスプレイ分数
\newcommand{\dpst}{\displaystyle}

\title{論文メモ
\large{\\--- 不動点定理と経済学？ ---}
}
\author{ナンブ　キトラ}
\date{\today}
\begin{document}
\maketitle

本棚を整理したいので，
溜まっている論文のメモを書いて未来への手紙とすることにする．


以下の論文の束は角谷の不動点定理を調べた時に
ついでに集めたものだと思う．
経済学の話なのでマジでよくわかんない．


シェルピンスキーの論文
\cite{Sierpinski1924}
もあるが，なぜこのラインナップに入っているのかは謎である．

\begin{thebibliography}{99}
\bibitem[和1]{1}
中嶋眞澄
「大数の強法則の初等的証明」, 
鹿児島経済論集, 45 (1) (2004), 1--5. 

\bibitem[和2]{2}
中嶋眞澄
「ゲーム理論におけるNashの定理の一般化」, 
鹿児島経済論集, 45 (3) (2004), 213--230. 

\bibitem[和3]{3}
飯村卓也
「離散不動点定理とその応用について」数理解析研究所講究録, 
1371(2004),  140--151. 
http://hdl.handle.net/2433/25481

\bibitem[和4]{4}
高橋渉
「不動点定理をめぐる最近の結果」, 
雑誌「数学」
28 (3) (1976), 
236--247, 
DOI: 10.11429/sugaku1947.28.236

\bibitem[和5]{5}
入谷純, 加茂知幸
「Arrow の不可能性定理と可換代数」，
http://www.f.waseda.jp/mtoda/Nikaido/pdf/algebraic.pdf

\end{thebibliography}



%READ!
%myplain is the same to amsplain style. 
\nocite{*}
\bibliographystyle{myplaindoidoi}
\bibliography{gameth.bib}



\end{document}